\documentclass[11pt]{article}
\usepackage[T1]{fontenc}
\usepackage[utf8]{inputenc}
\usepackage{lmodern,amsmath,amssymb,amsthm,microtype}
\usepackage[margin=25mm]{geometry}
\usepackage{hyperref}
\hypersetup{hidelinks,pdftitle={The finite-energy equation on the projective line over arbitrary complete fields}}
\setlength{\emergencystretch}{3em}
\newtheorem{theorem}{Theorem}
\newtheorem{lemma}[theorem]{Lemma}
\newtheorem{proposition}[theorem]{Proposition}
\newcommand{\MA}{\operatorname{MA}}
\newcommand{\PSH}{\operatorname{PSH}}
\newcommand{\E}{\mathcal E}
\newcommand{\an}{\mathrm{an}}
\title{The finite-energy equation on the projective line\\over arbitrary complete fields}
\author{P16.4 research continuation 222\\New partial result awaiting independent review}
\date{7 October 2026}
\begin{document}
\maketitle
\begin{abstract}
Using four complete independently accepted supporting theorems, we
construct the actual finite-energy potential of every finite-dual
Radon target on the projective line over any complete
characteristic-zero real-valued non-Archimedean field. Over a
nontrivially valued algebraically closed field, finite rooted trees
give decreasing continuous psh potentials with genuine prescribed
retracted measures. The dual directly bounds their energy costs.
An integration identity and radial upper semicontinuity prove
uniqueness at every analytic point, including terminal points.
Completed-algebraic-closure descent then handles arbitrary
nontrivial ground fields; the independently accepted trivial-field
theorem handles the trivial boundary. This proves the projective-line
sector only. No equation on a general variety is asserted.
\end{abstract}

\section{Exact prerequisites and statement}

The caller mathematical grants are the \emph{entire} accepted
manuscripts \cite{O,D,N,T}, with their source-preserving harmless
clarifications. Grant O supplies the genuine all-field actual
$\E^1$ operator, positivity and raw mass, integrability of all
$\E^1$ tests against its outputs, decreasing weighted continuity,
the energy identity, concavity, affine curvature calculus and
segment differentiation. Grant D supplies the completed algebraic
closure quotient, actual invariant psh and finite-energy descent,
the unique invariant Radon lift, and equality of continuous-test
duals. Grant N supplies the norming-reference target criterion and
establishes its hypotheses for standard projective space over
every complete field. Grant T supplies the actual ample equation
and every-solution uniqueness for all trivially valued
characteristic-zero fields. None of these grants supplies the new
nontrivial-field equation proved below.

Let $K$ be complete, of characteristic zero, with a real-valued
non-Archimedean absolute value. Initially take
$L=\mathcal O_{\mathbb P^1_K}(1)$, the standard coordinate FS
reference $\psi$, and $Z=(\mathbb P^1_K)^\an$. Write $\xi$ for the
Gauss point. Actual psh metrics are all-point decreasing FS-net
endpoints, excluding the identically $-\infty$ endpoint. A scalar
$u$ denotes $\psi+u$. The extended energy is the infimum of
continuous psh majorant energies; $\E^1$ means its actual finite
energy cone. By N,
\[
 \MA(0)=\delta_\xi,\qquad \sup_Zu=u(\xi)\in\mathbb R.
\]
We normalize by $u(\xi)=0$. For a positive Radon probability $\mu$
put
\[
 J(\mu)=\sup_{v\text{ continuous psh}}
                \left(E(v)-\int_Zv\,d\mu\right).
\]
All masses are raw masses, here equal to one.

\begin{theorem}\label{thm:main}
For every such $K$ and every positive Radon probability $\mu$ on
$Z$, the following are equivalent:
\begin{enumerate}
\item $J(\mu)<\infty$;
\item every actual $\E^1$ scalar is absolutely $\mu$-integrable;
\item $\mu=\MA(u)$ for an actual $u\in\E^1$.
\end{enumerate}
Every solution is unique modulo one finite real constant at
\emph{all} analytic points. The normalized solution maximizes
$E(v)-\int v\,d\mu$ on the actual $\E^1$ cone. The same assertions,
with target mass $\deg L$, hold for every actual ample rational
line bundle on $\mathbb P^1_K$ and every continuous psh reference.
\end{theorem}

The equivalence of the first two conditions is granted by N.
For every actual output, the second condition follows from O.
The work is actual existence from finite dual and every-solution
uniqueness. We do not assume compactness of an abstract energy
completion, continuity of envelopes or outer differentiation.

\section{Disk points and the rooted tree over the algebraic closure}

Until Section~\ref{sec:descent}, let $C$ be complete, algebraically
closed and nontrivially valued. Distances use the natural logarithm,
in the same normalization as metric weights. The elementary tree
description is recorded in \cite[Sections 1.3--1.6 and 2.2]{AWS}.
We reconstruct the properties used here. The earlier finite-tree
calculation in \cite[Sections 3.2--3.3]{BR} is compatible with them;
its original ground-field restriction is not used to extend a
theorem to arbitrary $C$.

For $a\in C$, $r>0$, the disk point $\zeta(a,r)$ has
\[
 |f|_{\zeta(a,r)}=\max_i|c_i|r^i,
 \qquad f(T)=\sum_i c_i(T-a)^i.
\]
The weighted Gauss norm is multiplicative. For radii in $|C^*|$,
scale to radius one and use reduction of polynomials; products of
nonzero residue polynomials stay nonzero. Other radii follow by
limits, since $|C^*|$ is dense in $\mathbb R_{>0}$: roots of an
element of nonunit norm give a dense divisible logarithmic subgroup.

Here is the disk classification needed at possible endpoints.
For any seminorm point $x$ on the affine chart, put
$r=\inf_{a\in C}|T-a|_x$. If the infimum is attained at $a$,
ultrametricity and minimality give
$|T-b|_x=\max(r,|a-b|)$ for every $b$; factorization of polynomials
then identifies $x$ with that disk point, allowing $r=0$.
If it is not attained, choose $a_j$ with
$r_j=|T-a_j|_x\downarrow r$. The disks $B(a_j,r_j)$ are nested
after choosing the numbers strictly decreasing. Indeed
$|a_i-a_j|\leq\max(r_i,r_j)$. For every $b$,
\[
 |T-b|_x\leq\max(r_j,|a_j-b|)
                \leq\max(r_j,|T-b|_x),
\]
so their disk seminorms tend to $x$ on all polynomials.
The disks have empty intersection, since a point in their
intersection would attain the infimum. Their limiting radius is
positive by completeness: if it were zero their centers would
form a Cauchy sequence with a point in every disk.
These are the terminal type IV points. Radius-zero disks are
the terminal type I points; adjoining infinity completes the
projective line.

Ordered chains of disks give arcs: between two disk points go to
the least containing disk and then down the other chain. Its radius
is $\max(r,s,|a-b|)$ for disks $(a,r),(b,s)$. Nested-disk endpoints
give the corresponding limits; the increasing-radius chain gives
the arc to infinity. Strict disk inequalities separate the branches,
so these arcs are unique. Comparable positive-radius disks have
length $|\log r-\log s|$. This constructs the rooted tree, without
asserting that its full path metric defines the analytic topology.
Every terminal point is approached on its root arc by finite
positive-radius disk points. Convergence follows by evaluating
linear polynomials in the displayed disk formulas, factorization,
and the second affine chart at infinity. Density of $|C^*|$ also
allows type II points on every nondegenerate finite subarc.

Write $g(x,y)$ for the length of
$[\xi,x]\cap[\xi,y]$, permitting infinity on a diagonal terminal
pair. For a finite disk point $y=\zeta(a,r)$ set
$b_y=\log\max(1,|a|,r)$ and
\begin{equation}\label{eq:kernel}
 u_y=\log\max\{|T_1-aT_0|,r|T_0|\}
              -\log\max\{|T_0|,|T_1|\}-b_y.
\end{equation}
This is a continuous relative FS scalar and satisfies
\begin{equation}\label{eq:atom}
                 u_y=-g(\,\cdot\,,y),\qquad
                 u_y(\xi)=0,\qquad \MA(u_y)=\delta_y.
\end{equation}
The first identity follows directly from disk radii. Inside the
unit disk its two sides are
$\log\max\{|T-a|,r\}$ minus their Gauss value. On the outer
branch the denominator contributes $\log\max(1,|T|)$.
If $|a|>\max(1,r)$, the path goes first from radius one to
$|a|$ and then down toward $a$, with total length
$2\log|a|-\log r$. These three disk configurations, and their
limits at endpoints and infinity, give the identity everywhere.

For the measure identity, first let $r=|c|$ with $c\in C^*$.
Use the standard $\mathbb P^1_{C^\circ}$ model with generic
coordinate sections $cT_0,T_1-aT_0$. Its generic identification is
an invertible linear change of homogeneous coordinates. The
model reference is the FS metric in \eqref{eq:kernel}, up to a
finite constant. Its reduced special fiber is a single projective
line, with $\mathcal O(1)$ degree one, and its generic Shilov point
projects to $\zeta(a,r)$. The actual model formula
\cite[Example 8.11 and (8.3)]{BE} therefore gives exactly
$\delta_y$. No field-degree or logarithmic-base factor enters.
For an arbitrary radius choose $r_j\in|C^*|$ tending to $r$.
Formula~\eqref{eq:kernel} converges uniformly, and the points
$\zeta(a,r_j)$ converge analytically to $y$. Uniform continuity
of the genuine continuous operator gives the same identity.

\section{Finite-tree potentials with genuine target measures}

Let $\Gamma$ range over convex hulls of $\xi$ and finitely many
finite disk points. These are compact finite metric trees; their
edges consist of positive-radius disk points. The family is
directed by inclusion. Let $r_\Gamma$ take $x$ to the last point
of $[\xi,x]$ in $\Gamma$. For every $y\in\Gamma$,
\begin{equation}\label{eq:retraction}
 g(r_\Gamma x,y)=g(x,y),\qquad
 g_\Gamma(x,z):=g(r_\Gamma x,r_\Gamma z).
\end{equation}
The retraction is continuous in the analytic topology. To check
this directly, take the finite leaves $y_1,\ldots,y_m$ of the
rooted tree. The functions $g(\,\cdot\,,y_i)=-u_{y_i}$ are
continuous and separate its points: for two comparable points
take a leaf beyond the farther one; for points on different
branches take a leaf beyond either branch. They embed the compact
tree into $\mathbb R^m$. The coordinates of $r_\Gamma x$ equal
$g(x,y_i)$, proving continuity. The kernel on $\Gamma^2$ is
continuous and obeys
\[
 |g(x,y)-g(x,z)|\leq\rho(y,z)\quad(x\in Z,\ y,z\in\Gamma).
\]
This follows because the shared initial root length changes by
at most the length between $y$ and $z$. Consequently $g_\Gamma$
is bounded and jointly continuous on $Z^2$.

For a Radon probability $\mu$, let
$\mu_\Gamma=(r_\Gamma)_*\mu$, a genuine positive Radon
probability. It can be diffuse along edges. Define
\[
 u_\Gamma(x)=-\int_\Gamma g(x,y)\,d\mu_\Gamma(y)
            =-\int_Zg_\Gamma(x,z)\,d\mu(z).
\]
Approximate this integral uniformly by finite convex sums of
the $u_y$: partition the finite tree into finitely many Borel
sets of small metric diameter, choose one point in each set,
and use its actual measure as coefficient. The displayed
Lipschitz bound controls the uniform error. Each finite sum is
continuous actual psh by real convexity on the fixed bundle.
Uniform closure gives a continuous actual psh $u_\Gamma$.
Affine curvature calculus and \eqref{eq:atom}, followed by
uniform continuity, give the \emph{genuine} identity
\begin{equation}\label{eq:GammaMA}
                     \MA(u_\Gamma)=\mu_\Gamma.
\end{equation}
The atomic approximations converge weakly to $\mu_\Gamma$ on
the compact finite tree; its inclusion into $Z$ is continuous.
This establishes the target measure rather than defining a
Hessian by naming it.

The measures $\mu_\Gamma$ tend weakly to $\mu$. Indeed the
continuous functions $g(\,\cdot\,,y)$ for all finite disk points
$y$ separate all points of $Z$. If two points lie on different
root branches, choose a disk point on one branch past their
meeting point; if they are comparable, choose a disk point
strictly beyond the nearer point. Terminal points are handled
by the disk approximations above. Together with constants,
the real algebra generated by these functions is uniformly
dense in $C(Z)$ by Stone--Weierstrass. For each element of this
algebra, retraction to every sufficiently large $\Gamma$ leaves
it unchanged by \eqref{eq:retraction}. Uniform approximation
then proves convergence against every continuous test. This is
a directed-net argument; no cofinal sequence is asserted.

\section{Actual existence from the continuous-test dual}

As $\Gamma$ increases, the common-root kernel $g_\Gamma$
increases: on each pair of arcs its value is the smaller of
their retained shared initial lengths. Hence $u_\Gamma$ decreases,
and $u_\Gamma(\xi)=0$. Its all-point endpoint $u$ is actual psh
by decreasing closure \cite[Corollary 7.5]{BE}, or its
reconstruction in N. The anchor excludes the identically
$-\infty$ endpoint.

Put
\[
 I_\Gamma=\int_{Z^2}g_\Gamma(x,z)\,d\mu(x)d\mu(z).
\]
The bounded continuous kernel justifies Fubini. The curve energy
identity and \eqref{eq:GammaMA} give
\begin{equation}\label{eq:cost}
 E(u_\Gamma)=-\tfrac12 I_\Gamma,\qquad
 \int_Zu_\Gamma\,d\mu=-I_\Gamma,\qquad
 E(u_\Gamma)-\int_Zu_\Gamma\,d\mu=\tfrac12 I_\Gamma.
\end{equation}
For the first equality use $u_\Gamma(\xi)=0$ and
$g_\Gamma(r_\Gamma x,r_\Gamma z)=g_\Gamma(x,z)$.
If $J(\mu)<\infty$, these continuous actual tests imply
$I_\Gamma\leq2J(\mu)$. Decreasing energy continuity therefore
gives $E(u)\geq-J(\mu)>-\infty$. Thus the endpoint is in the
\emph{actual} $\E^1$ cone. Grant O applies to its defining
decreasing continuous net and yields
\[
 \MA(u)=\lim_\Gamma\MA(u_\Gamma)
             =\lim_\Gamma\mu_\Gamma=\mu.
\]
No abstract Hilbert-space element is substituted for an actual
potential; no compactness or envelope theorem is used.

\section{Uniqueness at all points and the variational value}

For normalized actual curve potentials $p,q\in\E^1$, O gives
\begin{equation}\label{eq:symmetry}
                     \int_Zp\,\MA(q)=\int_Zq\,\MA(p).
\end{equation}
Here is a derivation from its exact granted statements.
$E(p)=\frac12\int p\,\MA(p)$ because $p(\xi)=0$.
Expand this identity on the actual segment $(1-t)p+tq$ using
affine curvature calculus. Its derivative at zero is
$-\int p\MA(p)+\frac12(\int p\MA(q)+\int q\MA(p))$.
The granted segment derivative is
$\int(q-p)\MA(p)$. Equality gives \eqref{eq:symmetry}.
Every integral is finite by O; common full-measure representatives
avoid subtraction of pole values.

If $\MA(p)=\mu$, apply \eqref{eq:symmetry} with $q=u_y$ for
any finite disk point. Since $\MA(u_y)=\delta_y$, O also
ensures $p(y)$ is finite, and
\begin{equation}\label{eq:values}
                         p(y)=-\int_Zg(z,y)\,d\mu(z).
\end{equation}
Thus every normalized solution has the same values at all finite
disk points. We must still recover its values at every terminal
point.

Every relative FS scalar is nonincreasing on each path away
from $\xi$. To verify this for a homogeneous section, factor
it over $C$ into linear forms. A nonzero linear form with zero
$a\in\mathbb P^1(C)$ has normalized log norm equal to its
Gauss log norm minus $g(x,a)$; this is the radius-zero limit
of \eqref{eq:kernel}, and for $a=\infty$ follows from the
second chart. Each common-root length increases along a
root path. Products add log norms, real weights add constants,
and finite maxima preserve nonincrease. This proves the claim
for every weighted FS system, and taking decreasing FS-net
endpoints proves it for every actual psh scalar.

Let $x$ be terminal and let finite disk ancestors $y_i$ tend
to $x$ on $[\xi,x]$. The preceding monotonicity gives
$p(y_i)\downarrow\ell\geq p(x)$ in the extended reals.
An actual psh scalar is upper semicontinuous, as an infimum
of continuous FS scalars, so
$p(x)\geq\limsup_i p(y_i)=\ell$. Consequently
$p(x)=\ell$, even when the value is $-\infty$.
Equation~\eqref{eq:values} determines all these limits.
The finite disk points together with terminal points exhaust
$Z$, proving all-point uniqueness modulo constants.

For any actual solution $u\in\E^1$, O supplies integrability
of every $\E^1$ test against $\mu=\MA(u)$. Concavity and
the granted derivative at the initial point of the segment
from $u$ to $v$ give
\[
 E(v)-E(u)\leq\int_Z(v-u)\,d\mu\qquad(v\in\E^1).
\]
Thus $u$ maximizes the stated functional. Continuous psh
majorants decreasing to $u$ converge in energy and in their
integrals against the fixed output $\mu$ by O's weighted
continuity, with constant curvature slots. Hence the
continuous-test supremum equals its value at $u$.
For the constructed solution, the same weighted continuity
and \eqref{eq:cost} give
\[
 J(\mu)=\tfrac12\sup_\Gamma I_\Gamma,
 \qquad E(u)=-\tfrac12\sup_\Gamma I_\Gamma,
 \qquad \int_Zu\,d\mu=-\sup_\Gamma I_\Gamma.
\]
These equalities involve finite quantities once the target
condition is established.

\section{Arbitrary ground fields and actual rational bundles}
\label{sec:descent}

Suppose first that $K$ is nontrivially valued and let
$C=\widehat{\overline K}$. Let
$\pi:(\mathbb P^1_C)^\an\to Z$ and
$G=\operatorname{Gal}(\overline K/K)$. All quotient and
invariance assertions in this paragraph are the full grant D,
including actual singular endpoints, rather than finite-level
or skeletal replacements. The unique invariant Radon lift
$\widetilde\mu$ has mass one and
\[
                  J_C(\widetilde\mu)=J_K(\mu)<\infty.
\]
The construction above gives an actual normalized upstairs
solution $\widetilde u\in\E_C^1$. The standard reference
and its Gauss point are $G$-invariant. For each $g\in G$,
naturality of the continuous operator and its defining
decreasing limit make $g^*\widetilde u$ a solution of the
same lifted target. Upstairs all-point uniqueness and the
common normalization imply $g^*\widetilde u=\widetilde u$.
Actual singular descent in D yields
$\widetilde u=\pi^*u$ with $u\in\E_K^1$. Genuine
covariance in O then gives $\MA(u)=\mu$.

If $v$ is any downstairs solution, its pullback is actual
finite energy by O. Its upstairs measure is invariant and
pushes to $\mu$; uniqueness of the invariant Radon lift in D
identifies it with $\widetilde\mu$. Upstairs uniqueness
therefore identifies the normalized pullback with
$\widetilde u$. Surjectivity of $\pi$ gives equality at
\emph{every} downstairs analytic point.
No averaging of singular potentials or of their pole values
has been performed. If $K$ is trivially valued, grant T
already establishes exactly this actual ample equation and
every-solution uniqueness; the nontrivial tree argument is
not applied to that case.

Finally $\operatorname{Pic}(\mathbb P^1_K)=\mathbb Z$. Indeed a
generic trivialization expresses a line bundle on this smooth curve
by a divisor. An affine closed point cut out by an irreducible
polynomial of degree $d$ is linearly equivalent to $d[\infty]$
by that polynomial's divisor. Hence every divisor is equivalent to
its degree times $[\infty]$, and a principal divisor has degree
zero. Consequently an
actual ample rational bundle has degree $a\in\mathbb Q_{>0}$
and is identified with $\mathcal O(a)$. Choose reference
$a\psi$. Clearing denominators in FS systems and scaling
their weights shows that $u=av$ identifies its actual psh
cone with that of $\mathcal O(1)$. Continuous curvature
linearity and the continuous energy identity give
\[
 \MA_a(av)=a\MA_1(v),\qquad E_a(av)=a^2E_1(v).
\]
The correspondence of continuous majorants extends the
energy identity to actual singular endpoints, and O's
decreasing continuity extends the operator identity.
For a raw-mass-$a$ target, set $\nu=\mu/a$; then
$J_a(\mu)=a^2J_1(\nu)$ and all-test integrability is
equivalent under the same scalar correspondence. Existence
and uniqueness follow. A change to any continuous psh
reference changes test scalars by a bounded continuous
function and energy by the reference cocycle in O.
The dual changes by a finite constant, so its finiteness,
the intrinsic condition, the absolute actual equation and
every-solution uniqueness are reference-independent.

\section{Boundary of this result}

Theorem~\ref{thm:main} concerns the actual ample projective-line
sector over all complete characteristic-zero real-valued
fields. It is not the full original P16.4 statement.
Higher-genus curves and arbitrary-dimensional varieties,
and the separate positive-residue-characteristic and
nonalgebraizable axes, are not closed here. The finite-tree
construction uses the single norming Gauss atom and
factorization of sections on the projective line; it is
not asserted for a general variety. All new mathematics in
this manuscript still requires independent review.

\begin{thebibliography}{99}\small
\bibitem{O} P16.4 continuation 217, \emph{An all-field finite-energy
mixed operator}, complete independently accepted manuscript,
review 179; normalization clarifications included.
\bibitem{D} P16.4 continuation 214, \emph{Completed-algebraic-closure
descent}, complete independently accepted manuscript, review 176.
\bibitem{N} P16.4 continuation 220, \emph{An intrinsic target
criterion from a norming atomic reference}, complete independently
accepted manuscript, review 182; source clarifications included.
\bibitem{T} P16.4 round 171, \emph{Galois descent at the
trivial-valuation boundary}, complete independently accepted
manuscript, review 173.
\bibitem{BE} S. Boucksom and D. Eriksson, \emph{Spaces of norms,
determinant of cohomology and Fekete points in non-Archimedean
geometry}, literal version 3,
\url{https://arxiv.org/pdf/1805.01016v3}.
\bibitem{AWS} M. Baker, \emph{An introduction to Berkovich analytic
spaces and non-archimedean potential theory on curves}, Arizona
Winter School chapter, author text,
\url{https://swc-math.github.io/aws/2007/BakerNotesMarch21.pdf}.
\bibitem{BR} R. Rumely and M. Baker, \emph{Analysis and dynamics
on the Berkovich projective line}, literal version 1,
\url{https://arxiv.org/pdf/math/0407433v1}.
\end{thebibliography}
\end{document}
